% ======================================================================
% CAPA PRINCIPAL (EXTERNA COM DESIGN GEOMÉTRICO)
% Arquivo: setup/pre_textual/capa_principal.tex
% ======================================================================
\thispagestyle{empty} % Remove numeração da página

% Altera temporariamente as margens para acomodar a barra lateral colorida
\newgeometry{top=1cm, bottom=2cm, left=5cm, right=0cm}

% --- BACKGROUND E GRAFISMOS (TIKZ) ---
\begin{tikzpicture}[remember picture,overlay]
    % Opções de separação
    \tikzset{x=10pt, y=10pt}
    
    % Loop para gerar as faixas coloridas na lateral esquerda (cores de 1 a 9)
    \foreach \i in {1,2,3,4,5,6,7,8,9}{
        \fill[color\i] 
        ([xshift={0.25*\paperwidth + (\i-1)*0.01*\paperwidth}]current page.north west) % Canto superior esquerdo deslocado
        rectangle 
        ([xshift={- 0.26*\paperwidth + (\i-1)*0.01*\paperwidth}]current page.south);    % Canto inferior direito deslocado
    }
    
    % Faixa sólida com a cor principal do eixo (color0)
    \fill [color0] (current page.north west) rectangle ([xshift = - 0.26\paperwidth]current page.south);

    % Inserção do ícone branco do eixo centralizado na barra lateral
    \node (anchor=center) at ([xshift = 0.125\paperwidth]current page.west) [opacity=1]{\includesvg[width=0.1\paperwidth]{\logoIconeEixoBranco}};
\end{tikzpicture}

% --- CONTEÚDO TEXTUAL E LOGOS DA CAPA ---
\begin{center}
    % Logo da UFPA no topo
    \includesvg[width = 0.1\linewidth]{\logoUfpa}
    \vspace{0.1\paperheight}
    
    % Logo texto do Nivelamento
    \includesvg[width = 0.5\linewidth]{\logoTextoNivelamento}
    \vspace{0.1\paperheight}

    % Nome do Eixo
    \LARGE
    \textbf{\eixo}
    \vspace{0.05\paperheight}

    % Assunto / Título da Apostila
    \begin{minipage}{0.6\paperwidth}
         \centering
         \assunto
    \end{minipage}
    
    \normalsize
    \vspace*{\fill} % Empurra as logos do rodapé para o limite inferior

    % --- RODAPÉ: PARCERIAS E APOIO ---
    \begin{minipage}{0.45\linewidth}
        \centering
        \includesvg[width = 0.3\linewidth]{\logoTutoria}
    \end{minipage}
    \hfill
    \begin{minipage}{0.45\linewidth}
        \centering
        \includesvg[width = 0.3\linewidth]{\logoItec}
    \end{minipage}
\end{center}

% Restaura as margens originais do documento para as páginas seguintes
\restoregeometry